\begin{theorem}[Theorem 2]
Let $n\ge2$, $p>1$ and $r>0$, and suppose the Sobolev--Poincar\'e inequality on balls holds in dimension $n$
(\noderef{CitedSobolevPoincareBall}).  Then there exist constants $C=C(p,r,n)>0$ and
$\alpha=\alpha(p,r,n)$ with $0<\alpha\le1$ with the following property.  Let
$\varepsilon\ge0$, $\gamma_0>0$, $\gamma_1\le r\gamma_0$, let
$\Omega\subseteq\R^n$ be open, let $A$ satisfy (Q1) of the thesis--(Q2) of the thesis on
$\Omega$ with parameters $p,\varepsilon,\gamma_0,\gamma_1$
(\noderef{EllipticStructureQ12}), let $u$ be a weak solution of
$\operatorname{div}A(\nabla u)=0$ in $\Omega$ (\noderef{IsWeakSolutionOn}) with
weak gradient $g=\nabla u$, let $\overline B_R(x_0)\subseteq\Omega$ and let
$0<\rho\le R$.  Then
\[
\operatorname*{osc}(\nabla u,B_\rho)\;\le\;C\,\sup_{B_R}|\nabla u|\,
\Bigl(\frac\rho R\Bigr)^{\alpha}.
\]
Both the oscillation and the supremum are essential (p.~6 of the paper), so the
statement is written as: for every $S$ with $|\nabla u|\le S$ a.e.\ in $B_R$,
\[
|g(x)-g(y)|\le C\,S\,(\rho/R)^{\alpha}
\quad\text{for a.e.\ }x\in B_\rho\text{ and a.e.\ }y\in B_\rho .
\]
Quantifying over all essential upper bounds $S$ is equivalent to using
$\esssup_{B_R}|\nabla u|$, which is itself such a bound.  The constants $C$ and
$\alpha$ are bound before $\varepsilon,\gamma_0,\gamma_1$ and all the data, and
the structure constants enter only through $\gamma_1\le r\gamma_0$: this is the
formal content of ``$C$ and $\alpha$ depend only on $p$, $\gamma_1/\gamma_0$ and
$n$''.

The hypothesis \noderef{CitedSobolevPoincareBall} is the cited result (C6) of the
paper's Note for the formalization, carried here as an explicit argument because
the proof uses it --- it enters through the constants of \noderef{OscillationDecayLemma} and \noderef{DegenerateAlternative} --- which is what the
Note permits for the results the thesis quotes without proof.  Although it is
bound before the existential for $C$ and $\alpha$, it introduces no dependence of
$C$ and $\alpha$ on any quantity outside the list the paper names: (C6) is a statement
about $n$ alone, and $n$ is already bound first, so the hypothesis carries no
numerical content of its own.

The cited solvability result (C2) (\noderef{CitedDirichletW1p}) is assumed here
for \emph{every} field $A'\colon\R^n\to\R^n$ and all parameters
$\varepsilon',\gamma_0',\gamma_1'$ \emph{with $\gamma_0'>0$ and $\gamma_1'>0$},
which is the form (C2) has in the paper's Note.  The hypothesis is bound before the constant, but it introduces no
dependence on any quantity outside the list the paper names: it mentions only
$n$ and $p$, both already bound, and it carries no numerical content of its own.
It is needed because the reduction of the case $\varepsilon=0$ to
$\varepsilon>0$ goes through \noderef{DegenerateLimitTransfer}, whose own (C2)
hypothesis is of that universal shape --- it has to be, since the regularized
fields to which (C2) is applied are produced inside that node's conclusion.

The cited regularity result (C1) (\noderef{CitedRegularityEpsPos}) is assumed in
the same way: for \emph{every} field $A'$, \emph{every} open set $\Omega'$ and
all parameters $\varepsilon',\gamma_0',\gamma_1'$ with
$\gamma_0'>0$ and $\gamma_1'>0$.  That is the form (C1) has in the
paper's Note, which states it for an arbitrary operator and an arbitrary
solution: ``If $\varepsilon>0$ and $A\colon\R^n\to\R^n$ satisfies
(Q1) of the thesis and (Q2) of the thesis, then every weak solution
$u\in W^{1,p}_{\mathrm{loc}}(\Omega)$ of (3.1) of the thesis satisfies
$u\in W^{2,2}_{\mathrm{loc}}(\Omega)$ and $\nabla u\in C(\Omega)$'' (Note for the
formalization, (C1)).  So the universal clause is the faithful transcription and
a clause attached only to the field $A$ and the domain $\Omega$ of this theorem
would be the weakening.  The clause is quantified over
$\varepsilon',\gamma_0',\gamma_1'$ with $\gamma_0'>0$, $\gamma_1'>0$ and over
open sets $\Omega'$ only, and those side conditions are not optional: the
structure conditions (Q1) of the thesis--(Q2) of the thesis carry positive constants in
the thesis (``$\gamma_0,\gamma_1>0$''), and (C1) is a statement about weak
solutions on an open set, so a clause quantified over $\gamma_0'\le0$ or over
arbitrary $\Omega'$ would assert more than the cited result does --- and since
this is a hypothesis of the present theorem, asserting more than the literature
gives risks making the theorem vacuously true.  The remaining antecedents need
no guard, because \noderef{CitedRegularityEpsPos} carries $\varepsilon>0$,
(Q1) of the thesis--(Q2) of the thesis and the weak-solution property inside its own
statement.  The same two positivity conditions are attached to the (C2) clause;
(C2) needs no openness condition, since \noderef{CitedDirichletW1p} is a
statement about balls and fixes its own domain.  Quantifying the set is not decorative: the proof consumes (C1) on
the balls $B_{R'}\subseteq B_R$ on which the alternative
\noderef{DegenerateAlternative} is applied --- that node asks for the predicate
itself at the ball --- and \noderef{CaccioppoliLevelWeight} takes the family
$D(\nabla u)$ as a parameter, which the proof must produce from the existential
inside (C1)'s conclusion.  Like (C2) and (C6), (C1) carries no numerical
content: it asserts membership and the existence of a continuous representative
of the gradient, and names no constant, so although
it is bound before the existentials for $C$ and $\alpha$ it introduces no
dependence of them on anything outside the list $p$, $\gamma_1/\gamma_0$, $n$
that the paper names.
\end{theorem}
\begin{proof}
Write $B_\sigma=B_\sigma(x_0)$ and $g^k=D_ku$.  Constants named $c$, $C$ or
$\gamma$ with a subscript are positive and depend only on $n$, $p$ and $r$;
this is checked at each occurrence.  $\omega_n=|B_1|$, so
$|B_\sigma|=\omega_n\sigma^{\,n}$ (Preamble: the volume of a ball).

The proof has two parts.  Part A fixes the constants; the exponent $\alpha$
among them is produced by \noderef{GradientOscillationEstimateEpsPos}, which is
the statement under the extra hypothesis $\varepsilon>0$ in a slightly more
flexible form (thesis pp.~27--29).  Part B removes $\varepsilon>0$.

\medskip
\noindent\emph{Part A: the constants.}

\textbf{(A0)} $S\ge0$.  The theorem's datum is a number $S$ with
$\|g\|\le S$ almost everywhere on $B_R$, and three steps below need $S$ to be
nonnegative.  It is: $|B_R|=\omega_nR^n>0$ (Preamble: the volume of a ball,
using $R\ge\rho>0$), so the set of $x\in B_R$ at which $\|g(x)\|\le S$ holds has
positive measure and is in particular nonempty; at any of its points
$0\le\|g(x)\|\le S$.

\textbf{(A1)} Apply \noderef{DegenerateLimitTransfer} with $n,p,r$ --- its
hypotheses are $2\le n$, $1<p$, $0<r$ --- and let $r'>0$ be the number it
produces; it depends only on $n,p,r$.  Put $r_+=\max\{r,r'\}>0$.  Every node
below that takes a ratio bound is instantiated at $r_+$, so that one set of
constants serves both the given problem, whose ratio bound is $r\le r_+$, and
the approximating problems of Part B, whose ratio bound is $r'\le r_+$.

\textbf{(A2)} Apply \noderef{GradientOscillationEstimateEpsPos} with $n,p,r_+$
--- its hypotheses there are $2\le n$, $1<p$ and $0<r_+$ --- together with the
Sobolev--Poincar\'e hypothesis \noderef{CitedSobolevPoincareBall} and the
universal cited regularity hypothesis (C1) (\noderef{CitedRegularityEpsPos}) of
the present theorem; those three are exactly its remaining head hypotheses, and
the last two are independent of the ratio bound, so they serve $r_+$ as well as
$r$.  Let $\alpha$ with $0<\alpha\le1$ be the exponent it produces; it depends
only on $n,p,r_+$, hence only on $n,p,r$.  Write (A2) also for its conclusion:
\emph{for every $\varepsilon>0$, $\gamma_0>0$ and $\gamma_1\le r_+\gamma_0$,
every open $\Omega_0$, every $A$ satisfying (Q1) of the thesis--(Q2) of the thesis on
$\Omega_0$ with parameters $p,\varepsilon,\gamma_0,\gamma_1$
(\noderef{EllipticStructureQ12}), every weak solution $u$ of
$\operatorname{div}A(\nabla u)=0$ in $\Omega_0$ (\noderef{IsWeakSolutionOn})
with weak gradient $g$, every $x_0$ and every $R_0>0$ with
$B_{R_0}(x_0)\subseteq\Omega_0$, every $T\ge0$ with $\|g\|\le T$ a.e.\ on
$B_{R_0}(x_0)$, every $0<\rho\le R_0$ and every $k$,}
\[
\operatorname*{osc}\bigl(g^k,B_\rho(x_0)\bigr)
\ \le\ 16\,T\,\Bigl(\frac\rho{R_0}\Bigr)^{\alpha},
\]
\emph{where $\operatorname*{osc}(f,E)=\operatorname{ess\,sup}_Ef-
\operatorname{ess\,inf}_Ef$ (\noderef{EssOscOn}).}  Two features of that
statement are used below: it asks only for the \emph{open} ball
$B_{R_0}(x_0)\subseteq\Omega_0$, which is what lets (B6$'$) apply it with
$\Omega_0$ a ball, and it does not require $\gamma_1>0$, that being derived
inside it from $\gamma_0>0$ and the structure conditions.

\textbf{(A3)} Apply \noderef{GradientSupEstimate} (Theorem~1) with $n,p,r_+$,
\noderef{CitedSobolevPoincareBall} and the two universal cited hypotheses (C2)
(\noderef{CitedDirichletW1p}) and (C1) (\noderef{CitedRegularityEpsPos}) of the
present theorem --- those three are exactly its remaining head hypotheses, and
they are independent of the ratio bound, so they serve $r_+$ as well as $r$ ---
and let $C_1>0$ be its constant.  Put
\[
\kappa=\bigl(C_1\,2^{\,2p-2}\,\omega_n\,2^{\,n}\bigr)^{1/p},
\qquad
C=\max\bigl\{16\sqrt n,\ 4\kappa\cdot16\sqrt n,\ 8\bigr\}
=\max\bigl\{64\sqrt n\,\kappa,\ 16\sqrt n,\ 8\bigr\} .
\]
All of $r_+,\alpha,C_1,\kappa,C$ are functions of $n$, $p$ and $r$ alone and
are fixed before $\varepsilon,\gamma_0,\gamma_1,S,\Omega,A,u,g,x_0,R,\rho$,
which is the quantifier order the statement asserts.  The order of choice is
acyclic: $r_+$ first, then $\alpha$ and $C_1$ at $r_+$, then $\kappa$ and
$C$.

\medskip
\noindent\emph{Part B: the general case $\varepsilon\ge0$.}

\textbf{(B1$'$) If $\varepsilon>0$.}  Apply (A2) with
$\Omega_0=\Omega$, $R_0=R$ and $T=S$; the inclusion $B_R\subseteq\Omega$ follows
from $\overline B_R\subseteq\Omega$, and $\gamma_1\le r\gamma_0\le r_+\gamma_0$.
Its hypothesis $T\ge0$ is (A0), its hypothesis $\|g\|\le T=S$ a.e.\ on $B_R$
is the theorem's datum, and $R_0=R>0$ because $R\ge\rho>0$.
For a.e.\ $x\in B_\rho$ and a.e.\ $y\in B_\rho$ one has
$\operatorname{ess\,inf}_{B_\rho}g^k\le g^k(x)\le
\operatorname{ess\,sup}_{B_\rho}g^k$ and the same for $y$, so
$|g^k(x)-g^k(y)|\le\operatorname*{osc}(g^k,B_\rho)$ for every $k$; hence
\[
\|g(x)-g(y)\|=\Bigl(\sum_k\bigl(g^k(x)-g^k(y)\bigr)^2\Bigr)^{1/2}
\le\sqrt n\max_k\operatorname*{osc}(g^k,B_\rho)
\le16\sqrt n\,S\Bigl(\frac\rho R\Bigr)^{\alpha}
\le C\,S\,\Bigl(\frac\rho R\Bigr)^{\alpha}.
\]

\textbf{(B2$'$) The normalisation $A(0)=0$.}  Put $\widetilde A(h)=A(h)-A(0)$.
Then $\widetilde A(0)=0$; $\widetilde A$ satisfies
(Q1) of the thesis--(Q2) of the thesis on $\Omega$ with the same
$p,\varepsilon,\gamma_0,\gamma_1$, because \texttt{hcont} and \texttt{hdiff} are
unaffected by subtracting a constant and \texttt{Q1} and \texttt{Q2} involve
only $\operatorname{fderiv}$ of the field, which is unchanged
(\noderef{EllipticStructureQ12}); and $u$ is a weak solution for
$\widetilde A$ in $\Omega$ iff it is one for $A$.  For the last point, the first
clause of \noderef{IsWeakSolutionOn} does not mention the field, and for a test
pair $(\varphi,g_\varphi)$ with $\varphi\in W^{1,p}(\Omega)$ of weak gradient
$g_\varphi$ (\noderef{MemW1pOn}), compact support and
$\operatorname{tsupport}\varphi\subseteq\Omega$, the two integrands differ by
$\sum_i(A(0))_ig_\varphi^{\,i}$ at every point.  That difference is integrable
on $\Omega$ (see below), so either both integrands are integrable --- in which
case the two integrals differ by the integral of the difference --- or neither
is, in which case both integrals are $0$ by the convention of the Bochner
integral and the two identities hold together.  In the first case the two
identities agree because
\[
\int_\Omega g_\varphi^{\,i}\,dx=0\qquad(i=1,\dots,n).
\]
Indeed $g_\varphi=0$ a.e.\ on the open set
$\Omega\setminus\operatorname{tsupport}\varphi$, where $\varphi$ vanishes
identically, because $0$ is also a weak gradient of $\varphi$ there and a weak
gradient is determined almost everywhere (Preamble: the fundamental lemma of the
calculus of variations, second consequence); in particular $g_\varphi^{\,i}$ is
integrable on $\Omega$, being supported in a compact set and in $L^p$ of it.
Now let $\delta_\flat>0$ be such that the closed
$\delta_\flat$-neighbourhood $K$ of $\operatorname{tsupport}\varphi$ is
contained in $\Omega$ --- such a $\delta_\flat$ exists because
$\operatorname{tsupport}\varphi$ is compact and $\Omega$ open --- and put
$\psi=\mathbf1_{K_{1/2}}*\rho_{\delta_\flat/4}$, where $K_{1/2}$ is the closed
$\delta_\flat/2$-neighbourhood of $\operatorname{tsupport}\varphi$ and
$\rho_{\delta_\flat/4}$ is the standard mollifier of the Preamble in dimension
$n$ at scale $\delta_\flat/4$.  Then $\psi$ is smooth (Preamble:
convolution with a smooth compactly supported kernel), $\psi=1$ on the closed
$\delta_\flat/4$-neighbourhood of $\operatorname{tsupport}\varphi$ and
$\psi=0$ off $K$, so $\psi$ is a legitimate test function for
\noderef{HasWeakGradientOn} on $\Omega$.  That identity, applied to $\psi$ and
the index $i$, reads
$\int_\Omega\varphi\,D_i\psi=-\int_\Omega g_\varphi^{\,i}\psi$.  Its left-hand
side vanishes, because $\psi$ is constant on a neighbourhood of
$\operatorname{tsupport}\varphi$ --- so $D_i\psi=0$ there --- while $\varphi=0$
off $\operatorname{tsupport}\varphi$; and its right-hand side is
$-\int_\Omega g_\varphi^{\,i}$, because $\psi=1$ on
$\operatorname{tsupport}\varphi$ and $g_\varphi^{\,i}=0$ a.e.\ off it.  This
proves the display.  Hence the two weak-solution properties coincide, and since
the conclusion of the theorem mentions only $g$, we may and do replace $A$ by
$\widetilde A$.  (This is thesis l.~542, ``Since $\dv A(0)=0$ we may assume
without loss of generality that $A(0)=0$''; the normalisation is internal to the
proof and is \emph{not} a hypothesis of the theorem.)

\textbf{(B3$'$) The structure on all of $\R^n$.}  $\Omega\ne\emptyset$, because
$\overline B_R\subseteq\Omega$ and $R\ge\rho>0$.  Of the six clauses of
\noderef{EllipticStructureQ12} for the field $(x,z,h)\mapsto\widetilde A(h)$,
the first two need no argument here: Carath\'eodory measurability in the point
holds because the field is constant in $x$, its value at $(x,z,h)$ being
$\widetilde A(h)$ with no occurrence of $x$, and a constant map is Borel
measurable; Carath\'eodory joint continuity holds because the field is constant
in $z$ as well, so $(z,h)\mapsto\widetilde A(h)$ is the composition of the
continuous projection $(z,h)\mapsto h$ with the continuous $\widetilde A$.
Neither mentions a domain, so both hold on $\R^n$ outright.  Each of the
remaining four clauses has the form ``for (almost) every $x$, if $x\in\Omega$ then $P$'' with $P$ a
statement that does not mention $x$ or $z$: \texttt{hcont} is
``$\widetilde A$ is continuous'', \texttt{hdiff} is
``$\widetilde A$ is $C^1$ on $\{\varepsilon+\|h\|^2\ne0\}$'', and \texttt{Q1},
\texttt{Q2} are inequalities in $h$ and $\xi$ alone.  The almost-everywhere
clauses hold at almost every $x$, and $|\Omega|\ge|B_R|>0$, so there is an $x\in\Omega$
at which they hold; the pointwise ones hold at every $x\in\Omega\ne\emptyset$.
Either way $P$ is true, hence holds for every $x\in\R^n$, and
$\widetilde A$ satisfies (Q1) of the thesis--(Q2) of the thesis on $\R^n=\mathrm{univ}$
with the same parameters.

\textbf{(B4$'$) The approximating sequence.}  Assume now $\varepsilon=0$.  By
clause (4) of \noderef{WeakGradientRestrict} with $q=p$ and the open sets
$\Omega\supseteq B_R$, together with the first clause of
\noderef{IsWeakSolutionOn} and the facts that $\overline B_R$ is compact and
$\overline B_R\subseteq\Omega$, we get $u\in W^{1,p}(B_R)$ with weak gradient
$g$ (\noderef{MemW1pOn}).  Apply \noderef{DegenerateLimitTransfer} --- the
instance of (A1) --- with $\gamma_0,\gamma_1$ (so $0<\gamma_0$ and
$\gamma_1\le r\gamma_0$), the field $\widetilde A$ with
$\widetilde A(0)=0$ by (B2$'$) and (Q1) of the thesis--(Q2) of the thesis at
$\varepsilon=0$ on $\mathrm{univ}$ by (B3$'$), the present theorem's universal
hypothesis (C2) (\noderef{CitedDirichletW1p}), the centre $x_0$, the radius
$R>0$, and the pair $(u,g)$ with $u\in W^{1,p}(B_R)$ and the weak-solution
property for $\widetilde A$ on $B_R$.  That last item is
\noderef{IsWeakSolutionRestrict}, applied with $n$, the exponent $p$, the open
set $\Omega$, the centre $x_0$, the radius $R$, the first field
$(x,z,h)\mapsto\widetilde A(h)$, the second field identically $0$ and the pair
$(u,g)$: its hypotheses are $1\le n$ and $1\le p$ (from $2\le n$ and $1<p$), the
openness of $\Omega$, $R>0$ (from $R\ge\rho>0$), the inclusion
$B_R\subseteq\Omega$ (from $B_R\subseteq\overline B_R\subseteq\Omega$), and the
weak-solution property for $\widetilde A$ on $\Omega$ --- which is the theorem's
datum for $A$ on $\Omega$ (\noderef{IsWeakSolutionOn}) transported to
$\widetilde A$ by (B2$'$).  It produces
$e_m>0$ with $e_m\to0$, fields $A_m$ with $A_m(0)=0$ and constants
$\gamma_0^{(m)}>0$, $\gamma_1^{(m)}\le r'\gamma_0^{(m)}\le r_+\gamma_0^{(m)}$
such that $A_m$ satisfies (Q1) of the thesis--(Q2) of the thesis on $\mathrm{univ}$ at
$e_m$, and functions $u_m$ with weak gradients $g_m$ that are weak solutions of
$\operatorname{div}A_m(\nabla u_m)=0$ in $B_R$
(\noderef{IsWeakSolutionOn}), with $u_m-u\in W^{1,p}_0(B_R)$ of weak gradient
$g_m-g$ (\noderef{MemW1pZeroOn}),
$\int_{B_R}\|g_m-g\|^p\,dx\to0$, and $g_m\to g$ a.e.\ on $B_R$.

\textbf{(B5$'$) A uniform gradient bound on $B_{R/4}$.}  Apply the conclusion of
\noderef{GradientSupEstimate} --- the instance of (A3), with the constant $C_1$
--- for each $m$, with $\varepsilon=e_m\ge0$, the constants
$\gamma_0^{(m)},\gamma_1^{(m)}$ (so $0<\gamma_0^{(m)}$ and
$\gamma_1^{(m)}\le r_+\gamma_0^{(m)}$), the open set $B_R$, the field $A_m$ with
$A_m(0)=0$, the solution $u_m$ with weak gradient $g_m$, the centre $x_0$ and
the radius $R/2>0$, whose closed ball $\overline B_{R/2}$ is contained in the
open $B_R$.  Its structure hypothesis on $B_R$ follows from the one on
$\mathrm{univ}$: each of the last four clauses of
\noderef{EllipticStructureQ12} is an assertion of the form ``for (almost) every
$x$, if $x\in\Omega'$ then \dots'', so replacing $\Omega'=\mathrm{univ}$ by the
subset $B_R$ only weakens the antecedent, while the first two clauses --- the
Carath\'eodory pair, i.e.\ measurability of $x\mapsto A_m(x,z,h)$ on $\R^n$ and
joint continuity of $(z,h)\mapsto A_m(x,z,h)$ at every $x\in\R^n$ --- carry no
domain guard and are literally the same assertion for $B_R$ as for
$\mathrm{univ}$.  Its conclusion is
\[
\bigl(e_m+\|g_m(x)\|^2\bigr)^{p/2}
\le\frac{C_1}{(R/2)^{\,n}}\int_{B_{R/2}}\bigl(e_m+\|g_m(y)\|^2\bigr)^{p/2}dy
\qquad\text{for a.e. }x\in B_{R/4}.
\]
Bound the integral.  Since $(e_m+t^2)^{p/2}\le(e_m^{1/2}+t)^{p}$ for $t\ge0$ and
$(a+b)^p\le2^{p-1}(a^p+b^p)$ for $a,b\ge0$ (convexity of $s\mapsto s^p$),
\[
\bigl(e_m+\|g_m\|^2\bigr)^{p/2}\le2^{p-1}\bigl(e_m^{p/2}+\|g_m\|^{p}\bigr)
\le2^{p-1}e_m^{p/2}+2^{2p-2}\bigl(\|g\|^{p}+\|g_m-g\|^{p}\bigr),
\]
so, writing $\varepsilon_m=\int_{B_R}\|g_m-g\|^p$ and using $\|g\|\le S$ a.e.\
on $B_R$ and $|B_{R/2}|\le|B_R|=\omega_nR^n$,
\[
\int_{B_{R/2}}\bigl(e_m+\|g_m\|^2\bigr)^{p/2}
\le2^{p-1}e_m^{p/2}\omega_nR^n+2^{2p-2}\omega_nR^nS^p+2^{2p-2}\varepsilon_m .
\]
Since $e_m\to0$ and $\varepsilon_m\to0$, the right-hand side has limit
$2^{2p-2}\omega_nR^nS^p$.  Also $\|g_m\|^p\le(e_m+\|g_m\|^2)^{p/2}$, so, putting
$S_m=\operatorname{ess\,sup}_{B_{R/4}}\|g_m\|$,
\[
S_m^{\,p}\le\frac{C_1\,2^{\,n}}{R^{\,n}}\int_{B_{R/2}}\bigl(e_m+\|g_m\|^2\bigr)^{p/2},
\qquad\text{hence}\qquad
\limsup_{m\to\infty}S_m\le\bigl(C_1\,2^{\,n}\,2^{2p-2}\omega_n\bigr)^{1/p}S=\kappa S .
\]
In particular $S_m<\infty$ for every $m$.  Also $S_m\ge0$: the function
$\|g_m\|$ is nonnegative everywhere and $|B_{R/4}|=\omega_n(R/4)^n>0$
(Preamble: the volume of a ball), so the essential supremum of a nonnegative
function over a set of positive measure is nonnegative --- it is an essential
upper bound for $\|g_m\|$ there, and an essential upper bound on a set of
positive measure is attained-or-exceeded at some point of it, where the function
is $\ge0$.  So $0\le S_m<\infty$, and $S_m$ is an admissible value of
the $T$ of (A2) for $\|g_m\|$ on $B_{R/4}$: both the
nonnegativity and the essential bound are in hand.

\textbf{(B6$'$) Passing (A2) to the limit.}  Let
$0<\rho\le R/4$.  For each $m$ apply (A2) with
$\varepsilon=e_m>0$, the constants $\gamma_0^{(m)},\gamma_1^{(m)}$, the open set
$\Omega_0=B_R$, the field $A_m$, the solution $u_m$ with weak gradient $g_m$,
the radius $R_0=R/4$ --- so $B_{R/4}\subseteq B_R$ --- and $T=S_m$.  It gives,
for every $k$, $\operatorname*{osc}(g_m^k,B_\rho)\le16S_m(4\rho/R)^{\alpha}$,
and hence, by the componentwise argument of (B1$'$), for a.e.\ $x$ and a.e.\
$y$ in $B_\rho$,
\[
\|g_m(x)-g_m(y)\|\le16\sqrt n\,S_m\Bigl(\frac{4\rho}R\Bigr)^{\alpha}
\le64\sqrt n\,S_m\Bigl(\frac\rho R\Bigr)^{\alpha},
\]
using $4^{\alpha}\le4$.  For each $m$ the set of $x\in B_\rho$ for which the
inner almost-everywhere statement fails is null, and for each such $x$ the set of
bad $y$ is null; intersecting over the countably many $m$ and with the
full-measure sets on which $g_m\to g$ (from (B4$'$)) leaves a full-measure set of
$x\in B_\rho$ and, for each of them, a full-measure set of $y\in B_\rho$, such
that the displayed bound holds for \emph{every} $m$ and $g_m(x)\to g(x)$,
$g_m(y)\to g(y)$.  For such $x,y$, passing to the limit in
$\|g_m(x)-g_m(y)\|\le64\sqrt n\,S_m(\rho/R)^{\alpha}$ and using (B5$'$),
\[
\|g(x)-g(y)\|=\lim_{m\to\infty}\|g_m(x)-g_m(y)\|
\le\liminf_{m\to\infty}64\sqrt n\,S_m\Bigl(\frac\rho R\Bigr)^{\alpha}
\le64\sqrt n\,\kappa\,S\Bigl(\frac\rho R\Bigr)^{\alpha}
\le C\,S\,\Bigl(\frac\rho R\Bigr)^{\alpha},
\]
the last inequality but one because
$\liminf_mS_m\le\limsup_mS_m\le\kappa S$ by (B5$'$).  For the remaining radii
$R/4<\rho\le R$ the trivial bound suffices: for a.e.\
$x,y\in B_\rho\subseteq B_R$ we have $\|g(x)\|\le S$ and $\|g(y)\|\le S$, so
$\|g(x)-g(y)\|\le2S$, while $(\rho/R)^{\alpha}>4^{-\alpha}\ge\tfrac14$ and
therefore $2S\le8S(\rho/R)^{\alpha}\le C\,S(\rho/R)^{\alpha}$ --- the first of
those two inequalities multiplies $\tfrac14<(\rho/R)^{\alpha}$ by $8S$ and so
needs $S\ge0$, which is (A0); for $S<0$ it would be false.  Together with
(B1$'$) this proves the theorem for every $\varepsilon\ge0$.
\end{proof}
